\clearpage
\section{확장탑과 비분리확장}
\label{sec:norm-trace-towers-inseparability}

\begin{learninggoal}
노름과 대각합이 확장탑을 따라 합성되는 이유를 이해한다. 분리차수와 비분리차수를 사용해 체매장 공식의 일반형을 얻고, 비분리확장에서 대각합이 왜 사라지는지 설명한다.
\end{learninggoal}

노름과 대각합은 기준체를 한 단계씩 내려도 일관되게 계산된다. 이 추이성은 복잡한 확장을 작은 확장들로 나누어 계산할 수 있게 해 주는 핵심 성질이다.

\subsection{확장탑의 추이성}

\begin{theorem}[노름과 대각합의 추이성]
\label{thm:norm-trace-transitivity}
유한 체확장의 탑
\[
  K\subseteq M\subseteq L
\]
이 주어졌다고 하자. 모든 $\alpha\in L$에 대해
\[
  \boxed{
  \Tr_{L/K}(\alpha)
  =\Tr_{M/K}\bigl(\Tr_{L/M}(\alpha)\bigr)}
\]
이고
\[
  \boxed{
  \operatorname N_{L/K}(\alpha)
  =\operatorname N_{M/K}\bigl(\operatorname N_{L/M}(\alpha)\bigr).}
\]
\end{theorem}

\begin{proofidea}
먼저 모든 확장이 분리인 경우를 생각한다. $L$의 $K$-매장들을 $M$에 대한 제한값에 따라 묶으면, 각 묶음의 합과 곱이 각각 $\Tr_{L/M}(\alpha)$와 $\operatorname N_{L/M}(\alpha)$의 켤레가 된다. 일반 경우에는 비분리차수가 만드는 공통 중복도를 함께 세면 같은 공식이 성립한다.
\end{proofidea}

\begin{proof}
먼저 $L/K$가 분리라고 하자. $M$의 각 $K$-매장 $\tau:M\to\overline K$는 $L$의 $K$-매장으로 정확히 $[L:M]$가지 방식으로 연장된다. $L$의 모든 $K$-매장을 $M$에 대한 제한 $\tau$에 따라 묶으면
\[
  \sum_{\substack{\sigma:L\to\overline K\\\sigma|_M=\tau}}
  \sigma(\alpha)
  =\tau\bigl(\Tr_{L/M}(\alpha)\bigr)
\]
이고
\[
  \prod_{\substack{\sigma:L\to\overline K\\\sigma|_M=\tau}}
  \sigma(\alpha)
  =\tau\bigl(\operatorname N_{L/M}(\alpha)\bigr)
\]
이다. 모든 $\tau$에 대해 다시 합하거나 곱하고 정리 \ref{thm:norm-trace-embeddings}를 적용하면 두 추이식을 얻는다.

일반 유한확장에서는 아래 정리 \ref{thm:norm-trace-inseparable-embeddings}의 공식에 비분리차수를 곱해 사용한다. 분리차수와 비분리차수는 확장탑에서 각각 곱셈적이므로, 위의 묶음 계산에 나타나는 중복도도 양변에서 일치한다. 따라서 같은 두 공식이 임의의 유한확장에 대해 성립한다.
\end{proof}

\begin{example}[이중 이차확장에서 단계별 계산]
$L=\Q(\sqrt2,\sqrt3)$, $M=\Q(\sqrt2)$, $\alpha=\sqrt2+\sqrt3$라 하자. $L/M$의 비자명한 자동동형은 $\sqrt3$의 부호를 바꾸므로
\[
  \Tr_{L/M}(\alpha)=2\sqrt2,
  \qquad
  \operatorname N_{L/M}(\alpha)
  =(\sqrt2+\sqrt3)(\sqrt2-\sqrt3)=-1.
\]
이제 $M/\Q$에 대해
\[
  \Tr_{M/\Q}(2\sqrt2)=0,
  \qquad
  \operatorname N_{M/\Q}(-1)=1.
\]
따라서
\[
  \Tr_{L/\Q}(\alpha)=0,
  \qquad
  \operatorname N_{L/\Q}(\alpha)=1.
\]
\end{example}

\begin{corollary}
$\alpha\in M$이면
\[
  \Tr_{L/M}(\alpha)=[L:M]\alpha,
  \qquad
  \operatorname N_{L/M}(\alpha)=\alpha^{[L:M]}.
\]
따라서
\[
  \Tr_{L/K}(\alpha)
  =[L:M]\Tr_{M/K}(\alpha),
\]
\[
  \operatorname N_{L/K}(\alpha)
  =\operatorname N_{M/K}(\alpha)^{[L:M]}.
\]
\end{corollary}

\subsection{비분리차수가 만드는 중복도}

유한확장 $L/K$의 분리차수와 비분리차수를 각각
\[
  [L:K]_{\mathrm s},
  \qquad
  [L:K]_{\mathrm i}
\]
로 쓰면
\[
  [L:K]=[L:K]_{\mathrm s}[L:K]_{\mathrm i}
\]
이다. $K$-매장의 개수는 $[L:K]_{\mathrm s}$이다.

\begin{theorem}[비분리확장을 포함한 매장 공식]
\label{thm:norm-trace-inseparable-embeddings}
$L/K$가 유한확장이고
\[
  q=[L:K]_{\mathrm i},
  \qquad
  \sigma_1,\dots,\sigma_s:L\longrightarrow\overline K
\]
가 서로 다른 모든 $K$-매장이라고 하자. 그러면 모든 $\alpha\in L$에 대해
\[
  \boxed{
  P_{\alpha,L/K}(T)
  =\prod_{j=1}^{s}
  \bigl(T-\sigma_j(\alpha)\bigr)^q.}
\]
따라서
\[
  \boxed{
  \Tr_{L/K}(\alpha)
  =q\sum_{j=1}^{s}\sigma_j(\alpha),
  \qquad
  \operatorname N_{L/K}(\alpha)
  =\left(\prod_{j=1}^{s}\sigma_j(\alpha)\right)^q.}
\]
\end{theorem}

\begin{proof}
$E=K(\alpha)$라 하자. $\alpha$의 최소다항식은 대수적 폐포에서 서로 다른 근마다 $[E:K]_{\mathrm i}$의 중복도를 갖는다. 한편 $E$의 각 $K$-매장은 $L$의 $K$-매장으로 $[L:E]_{\mathrm s}$가지 방식으로 연장된다. 따라서
\[
  \prod_{j=1}^{s}(T-\sigma_j(\alpha))
\]
에서는 최소다항식의 각 서로 다른 근이 $[L:E]_{\mathrm s}$번 나타난다.

이 곱을 $q=[L:K]_{\mathrm i}
=[E:K]_{\mathrm i}[L:E]_{\mathrm i}$제곱하면 각 근의 중복도는
\[
  [E:K]_{\mathrm i}[L:E]_{\mathrm s}[L:E]_{\mathrm i}
  =[E:K]_{\mathrm i}[L:E]
\]
가 된다. 이는
\[
  f_\alpha(T)^{[L:E]}
  =P_{\alpha,L/K}(T)
\]
에서의 중복도와 같다. 특성다항식의 둘째 계수와 상수항을 비교하면 결론을 얻는다.
\end{proof}

\begin{corollary}[비분리확장의 대각합]
\label{cor:inseparable-trace-zero}
$L/K$가 유한 비분리확장이면
\[
  \Tr_{L/K}=0
\]
인 영사상이다.
\end{corollary}

\begin{proof}
$\charac K=p>0$이고 비분리차수 $q$는 $p$의 양의 거듭제곱이므로 $q=0$ in $K$이다. 정리 \ref{thm:norm-trace-inseparable-embeddings}의 대각합 공식에 적용한다.
\end{proof}

\begin{corollary}[순수비분리확장]
\label{cor:purely-inseparable-norm-trace}
$L/K$가 차수 $q=p^r>1$의 유한 순수비분리확장이면 모든 $\alpha\in L$에 대해
\[
  \Tr_{L/K}(\alpha)=0,
  \qquad
  \operatorname N_{L/K}(\alpha)=\alpha^q.
\]
특히 $\alpha^q\in K$이다.
\end{corollary}

\begin{proof}
순수비분리확장은 대수적 폐포로 가는 $K$-매장이 하나뿐이며 그것은 포함사상이다. 앞 정리에 $s=1$을 대입한다.
\end{proof}

\begin{example}
$K=\F_p(t)$이고 $L=K(u)$, $u^p=t$라 하자. $L/K$는 차수 $p$의 순수비분리확장이다. 따라서
\[
  \Tr_{L/K}(u)=0,
  \qquad
  \operatorname N_{L/K}(u)=u^p=t.
\]
최소다항식 $T^p-t$은 대수적 폐포에서 $(T-u)^p$이므로 특성다항식의 모든 근이 한 점에 중복되어 있다.
\end{example}

\begin{pitfall}
$\Tr_{L/K}(1)=0$이라는 사실만으로 $L/K$가 비분리라고 결론내릴 수는 없다. 분리확장에서도 $\charac K$가 $[L:K]$를 나누면 $\Tr(1)=0$일 수 있다. 비분리확장의 특징은 \emph{모든 원소의 대각합이} $0$이라는 것이다.
\end{pitfall}

\begin{keyidea}
분리확장에서는 켤레들이 서로 다른 좌표로 나타난다. 비분리확장에서는 같은 좌표가 $p$의 거듭제곱만큼 중복되고, 그 중복도가 표수에서 $0$이 되어 대각합 전체를 지운다.
\[
  \boxed{
  L/K\text{ finite inseparable}
  \quad\Longrightarrow\quad
  \Tr_{L/K}=0.}
\]
\end{keyidea}

\begin{checkpoint}
\begin{enumerate}[label=(\alph*)]
  \item $K\subseteq M\subseteq L$에서 $\alpha\in M$일 때 두 추이식이 위 따름정리의 공식으로 줄어드는 과정을 적어라.
  \item $K=\F_p(s,t)$, $L=K(s^{1/p},t^{1/p})$일 때 $[L:K]$, $\Tr_{L/K}(s^{1/p}+t^{1/p})$, $\operatorname N_{L/K}(s^{1/p}t^{1/p})$를 구하라.
  \item 분리확장 $\F_{p^p}/\F_p$에서 $\Tr(1)=0$이지만 대각합 사상이 영이 아닌 이유를 설명하라.
\end{enumerate}
\end{checkpoint}
